% ------------------------------------------------------------
% Compile with XeLaTeX (system fonts are used).
%   tectonic proposal.tex      or      xelatex proposal.tex
% ------------------------------------------------------------
% !TeX program = xelatex
% !TeX root = proposal.tex
%
% Wits adaptation of the BITM LaTeX Proposal Template v4
% (github.com/eduankotze/BITM), which is a UFS Department of Computer
% Science & Informatics template. Section structure and page budgets are
% kept from that template. Branding is rebuilt for Wits.
%
% NOTE ON SECTION 9. The source template frames methodology through
% Saunders' research onion and CRISP-DM, which suit applied information
% systems and data science work. This proposal proves theorems and then
% tests them, so 9.7 is restructured around the actual work phases.
% The deviation is deliberate and is flagged in that section.

\LoadClass{cover_and_title_pages}
\usepackage{proposal_template_style}
\usepackage{changepage}
\usepackage[none]{hyphenat}

\begin{document}

	% === Cover page fields ===
	\renewcommand{\proposal}{PhD Proposal}
	\renewcommand{\schoolname}{School of Statistics and Actuarial Science}
	\renewcommand{\facultyname}{Faculty of Science}
	\renewcommand{\thesistitle}{Pricing and Hedging through Benchmark Transitions in the Absence of Liquid Option Markets}
	\renewcommand{\initials}{T.B.J.}
	\renewcommand{\surname}{Baloyi}
	\renewcommand{\studentnum}{1113941}
	\renewcommand{\submissiondate}{September 2026}

	% === Title page fields ===
	\renewcommand{\fullnames}{Baloyi Thabang Bongani Junior}
	\renewcommand{\degree}{PHILOSOPHIAE DOCTOR}
	\renewcommand{\proposaltype}{(Mathematical Statistics)}
	\renewcommand{\universityname}{University of the Witwatersrand}
	\renewcommand{\campusname}{Johannesburg, South Africa}
	\renewcommand{\supervisor}{\hl{To be confirmed}}
	\renewcommand{\cosupervisor}{\hl{To be confirmed}}
	\renewcommand{\proposalmonth}{September}
	\renewcommand{\proposalyear}{2026}

	\maketitlepage

	\pagenumbering{roman}

	\renewcommand{\contentsname}{Table of Contents}
	\tableofcontents
	\newpage

	\renewcommand{\listfigurename}{List of Figures}
	\listoffigures

	\renewcommand{\listtablename}{List of Tables}
	\listoftables
	\newpage

	\section*{List of Abbreviations}
	\addcontentsline{toc}{section}{List of Abbreviations}

	\begin{tabular}{@{}ll@{}}
		CAS      & Credit adjustment spread \\
		ISDA     & International Swaps and Derivatives Association \\
		JIBAR    & Johannesburg Interbank Average Rate \\
		LIBOR    & London Interbank Offered Rate \\
		MPG      & Market Practitioners Group \\
		OIS      & Overnight index swap \\
		RFR      & Risk-free rate \\
		SABR     & Stochastic alpha, beta, rho \\
		SARB     & South African Reserve Bank \\
		SOFR     & Secured Overnight Financing Rate \\
		ZARONIA  & South African Rand Overnight Index Average \\
	\end{tabular}

	\newpage

	\section*{Glossary}
	\addcontentsline{toc}{section}{Glossary}

	\begin{description}
		\item[Backward-looking rate] A rate compounded from overnight fixings over an accrual period, known only at the end of that period.
		\item[Compatible model set] The set of admissible models whose implied observations agree with the market observations within stated quote tolerances.
		\item[Identified functional] A target valuation that takes the same value at every model in the compatible set.
		\item[Identification range] The interval between the smallest and largest target values over the compatible set. Not a confidence interval.
		\item[Hedging loss] The discounted liability less discounted terminal wealth of the hedge; positive values are shortfalls.
		\item[Marginal coverage] A probability averaged over calibration samples, as distinct from coverage conditional on the realised calibration sample.
	\end{description}

	\newpage
	\pagenumbering{arabic}

	% ============================================================
	\section{Introduction} % Approx. 1 page
	\todo[inline]{Draw from the concept note: why replacing a benchmark changes contractual cashflows, and why the successor option market may be absent or thin.}

	% ============================================================
	\section{Background} % Approx. 2-3 pages
	\todo[inline]{Benchmark reform and the JIBAR--ZARONIA transition; SARB and MPG conventions; what curves, swaps, basis instruments and legacy options reveal; interest-rate modelling approaches; conformal prediction under dependence.}

	% ============================================================
	\section{Research Gaps} % Approx. 0.5 page
	\todo[inline]{The provisional gap: connecting a rigorously characterised compatible model set to a finite-sample analysis of the hedging losses those models imply. State that the gap is provisional pending the structured literature comparison.}

	% ============================================================
	\section{Research Problem and Statement} % Approx. 0.75 - 1 page
	\todo[inline]{A fitted model returns one successor price although the observations are compatible with several; small average hedging error does not imply a valid upper-tail guarantee.}

	% ============================================================
	\section{Thesis Statement (PhD only)} % Approx. 1-2 sentences
	\todo[inline]{One or two sentences. Keep it falsifiable and do not assert what is not yet proved.}

	% ============================================================
	\section{Research Aim} % Approx. 0.5 page

	To develop and evaluate a framework for identifying successor-rate option
	valuations and quantifying residual hedging risk under limited
	option-market information.

	% ============================================================
	\section{Research Questions and Objectives} % Approx. 1-2 pages

	\textbf{RQ1.} Which successor valuation functionals are constant across the
	models compatible with legacy option quotes and curve or basis observations?

	\textbf{RQ2.} Which additional observations or restrictions restore
	identification, and how stable is that conclusion to quote noise?

	\textbf{RQ3.} Under what dependence and distribution-shift assumptions can a
	calibrated loss threshold retain a finite-sample guarantee?

	\textbf{RQ4.} Are the resulting price ranges and risk bounds informative for
	a realistic hedge, once costs and data limitations are included?

	\medskip
	\textbf{Objective 1.} Specify a coherent transition model, contractual
	conventions and observable market information, keeping risk-neutral pricing
	assumptions separate from historical estimation.

	\textbf{Objective 2.} Prove identification or non-identification results for
	a restricted model class, and determine the resulting valuation ranges or
	approximation errors.

	\textbf{Objective 3.} Construct implementable hedges from available
	instruments and investigate finite-sample loss guarantees under explicit
	assumptions.

	\textbf{Objective 4.} Evaluate pricing uncertainty, hedge performance and
	certificate informativeness in simulation and in an appropriately limited
	empirical application.

	% ============================================================
	\section{Hypotheses (Optional)} % Approx. 0.5 page
	\todo[inline]{Optional. If included, state them so that they can fail, and label them as conjectures rather than results.}

	% ============================================================
	\section{Research Design and Methodology} % Approx 4-5 pages

	\subsection{Introduction}
	\todo[inline]{The source template introduces Saunders' research onion here. Retain it only if the supervisor expects that framing; a mathematical-statistics proposal may instead state the design directly.}

	\subsection{Research Philosophy}
	\todo[inline]{State the position honestly rather than selecting a label for its own sake.}

	\subsection{Approach to Theory Development}
	Deductive. Results are derived from stated assumptions and then tested
	numerically and, where data permit, empirically. Simulation supports a
	proof; it does not replace one.

	\subsection{Methodological Choice}
	Quantitative throughout: mathematical analysis, controlled simulation and
	an empirical case study on market data.

	\subsection{Research Strategy}
	Mathematical derivation of identification and coverage results, followed by
	a simulation study with known data-generating laws, followed by an empirical
	application limited to what the data audit supports.

	\subsection{Time Horizon}
	Longitudinal. The loss process is observed chronologically, and training,
	method selection, calibration and evaluation periods are separated in time.

	\subsection{Techniques and Procedures}

	The source template organises this subsection around CRISP-DM, whose phases
	suit an applied data-mining project. The work proposed here proves results
	and then tests them, so the phases below replace that cycle. The substitution
	is deliberate.

	\subsubsection{Model and Observation Map}
	\todo[inline]{Low-dimensional Markov model in an overnight-rate factor and a spread factor under an explicit pricing measure; the observation map must price actual cashflows before any identification claim.}

	\subsubsection{Identification Analysis}
	\todo[inline]{Compatible set, identified functionals, identification ranges; constructive non-identification examples preserving the full observation map.}

	\subsubsection{Hedging and the Loss Process}
	\todo[inline]{Admissible strategies, discounted gains, transaction costs, the signed loss definition; sensitivity and mean--variance benchmarks.}

	\subsubsection{Finite-Sample Guarantees}
	\todo[inline]{Split conformal baseline, the Beta law for calibration-conditional coverage, then the dependence and shift corrections.}

	\subsubsection{Simulation Study}
	\todo[inline]{Known data-generating laws; dependence and shift varied separately then together; ties and overlapping horizons included.}

	\subsubsection{Empirical Evaluation}
	\todo[inline]{Data audit, chronological separation, locked test period; report what was actually validated.}

	% ============================================================
	\section{Validity and Reliability} % Approx. 0.5 page

	\subsection{Validity}
	\todo[inline]{Whether the constructions answer the stated questions: full observation maps in counterexamples, correct conditioning in coverage statements.}

	\subsection{Reliability}
	\todo[inline]{Reproducibility: versioned code, configurations, seeds, Monte Carlo error estimates, documented data cleaning.}

	% ============================================================
	\section{Ethical Considerations} % Approx. 0.5 page

	The research uses market data and simulation rather than human participants,
	so no participant consent arises. The applicable institutional ethics or
	exemption process will still be confirmed. Data licences, attribution and the
	treatment of confidential market information will be documented. AI
	assistance in preparing this proposal and the supporting manuscript is
	disclosed as the University requires; responsibility for every retained claim
	rests with the candidate.

	% ============================================================
	\section{Expected Contributions of the Research} % Approx. 0.5 page

	\subsection{Theoretical Contributions}
	\todo[inline]{Identification or non-identification for a stated model class and observation set; a justified finite-sample loss guarantee, or a credited application of existing theory. Label as proposed until proved.}

	\subsection{Practical Contributions}
	\todo[inline]{What a valuation range and a loss threshold would let a practitioner say, and what they would not.}

	% ============================================================
	\section{Limitations of the Study} % Approx. 0.5 page
	\todo[inline]{Restricted model class, data access, dependence assumptions, unobserved successor prices, transfer across currencies.}

	% ============================================================
	\section{Timeline} % Approx. 0.5 page

	\begin{table}[H]
		\centering
		\caption{Indicative full-time schedule.}
		\label{tab:timeline}
		\begin{tabular}{@{}ll@{}}
			\hline
			\textbf{Months} & \textbf{Work} \\
			\hline
			1--6   & Foundations, novelty review and data audit \\
			7--14  & Transition model and identification \\
			15--22 & Hedging and statistical analysis \\
			23--29 & Simulation and empirical evaluation \\
			30--36 & Synthesis, revision and examination preparation \\
			\hline
		\end{tabular}
	\end{table}

	Month 6 is the first decision point: confirm data access, the restricted
	model, and at least one candidate original result before expanding
	implementation.

	% ============================================================
	\section{Budget} % Approx. 0.5 page
	\todo[inline]{Data licensing and computing costs to be priced before the final budget is submitted. No funding or vendor access is assumed.}

	% ============================================================
	\section{Computational Resources} % Approx. 0.5 page
	\todo[inline]{Scientific computing software and a workstation; state whether cluster access is needed for the simulation study.}

	% ============================================================
	\section{Chapter Layout} % Approx. 1 page
	\todo[inline]{The ten-chapter outline. State that the chapter count is an organisational choice and not a claim about doctoral depth.}

	% ============================================================
	\newpage
	\bibliographystyle{apacite}
	\bibliography{library}
	\addcontentsline{toc}{section}{References}

\end{document}
